% !TeX root = closed-systems-from-comparison-completeness.tex
% ============================================================
\chapter*{Orientation to Part VII}
\phantomsection
\label{part:intro-numerical}
% ============================================================
\begin{chapterorientation}
\orientationitem{Settled}{The electromagnetic and matter-sector chapters provide charge, scalar, and refinement-depth data.}
\orientationitem{Task}{Part~VII develops the conditional numerical closure layer and separates intrinsic closure from realization-level normalization.}
\orientationitem{Handoff}{The output is the scalar closure chain, the same-channel correction structure, and the operator-realization boundary.}
\end{chapterorientation}

Part~VII is written with its status distinctions explicit.  It records
what the closed transport stack fixes internally and what remains a
separate realization or boundary datum.  The numerical layer is therefore
not treated as an external fit to the preceding theory.

\Cref{ch:matter-numbers} develops the quartic carrier and mass-depth
law.  \Cref{ch:numerical-closure} organizes the unified scalar channel
and same-channel correction structure.  \Cref{ch:matter-realization}
then states the operator-first realization problem and the conditions
under which the conditional package would become an operator-level
result.
